% Options for packages loaded elsewhere
\PassOptionsToPackage{unicode}{hyperref}
\PassOptionsToPackage{hyphens}{url}
\documentclass[
]{article}
\usepackage{xcolor}
\usepackage{amsmath,amssymb}
\setcounter{secnumdepth}{-\maxdimen} % remove section numbering
\usepackage{iftex}
\ifPDFTeX
  \usepackage[T1]{fontenc}
  \usepackage[utf8]{inputenc}
  \usepackage{textcomp} % provide euro and other symbols
\else % if luatex or xetex
  \usepackage{unicode-math} % this also loads fontspec
  \defaultfontfeatures{Scale=MatchLowercase}
  \defaultfontfeatures[\rmfamily]{Ligatures=TeX,Scale=1}
\fi
\usepackage{lmodern}
\ifPDFTeX\else
  % xetex/luatex font selection
\fi
% Use upquote if available, for straight quotes in verbatim environments
\IfFileExists{upquote.sty}{\usepackage{upquote}}{}
\IfFileExists{microtype.sty}{% use microtype if available
  \usepackage[]{microtype}
  \UseMicrotypeSet[protrusion]{basicmath} % disable protrusion for tt fonts
}{}
\makeatletter
\@ifundefined{KOMAClassName}{% if non-KOMA class
  \IfFileExists{parskip.sty}{%
    \usepackage{parskip}
  }{% else
    \setlength{\parindent}{0pt}
    \setlength{\parskip}{6pt plus 2pt minus 1pt}}
}{% if KOMA class
  \KOMAoptions{parskip=half}}
\makeatother
% definitions for citeproc citations
\NewDocumentCommand\citeproctext{}{}
\NewDocumentCommand\citeproc{mm}{%
  \begingroup\def\citeproctext{#2}\cite{#1}\endgroup}
\makeatletter
 % allow citations to break across lines
 \let\@cite@ofmt\@firstofone
 % avoid brackets around text for \cite:
 \def\@biblabel#1{}
 \def\@cite#1#2{{#1\if@tempswa , #2\fi}}
\makeatother
\newlength{\cslhangindent}
\setlength{\cslhangindent}{1.5em}
\newlength{\csllabelwidth}
\setlength{\csllabelwidth}{3em}
\newenvironment{CSLReferences}[2] % #1 hanging-indent, #2 entry-spacing
 {\begin{list}{}{%
  \setlength{\itemindent}{0pt}
  \setlength{\leftmargin}{0pt}
  \setlength{\parsep}{0pt}
  % turn on hanging indent if param 1 is 1
  \ifodd #1
   \setlength{\leftmargin}{\cslhangindent}
   \setlength{\itemindent}{-1\cslhangindent}
  \fi
  % set entry spacing
  \setlength{\itemsep}{#2\baselineskip}}}
 {\end{list}}
\usepackage{calc}
\newcommand{\CSLBlock}[1]{\hfill\break\parbox[t]{\linewidth}{\strut\ignorespaces#1\strut}}
\newcommand{\CSLLeftMargin}[1]{\parbox[t]{\csllabelwidth}{\strut#1\strut}}
\newcommand{\CSLRightInline}[1]{\parbox[t]{\linewidth - \csllabelwidth}{\strut#1\strut}}
\newcommand{\CSLIndent}[1]{\hspace{\cslhangindent}#1}
\setlength{\emergencystretch}{3em} % prevent overfull lines
\providecommand{\tightlist}{%
  \setlength{\itemsep}{0pt}\setlength{\parskip}{0pt}}
\newcommand\R{\mathbb{R}} \newcommand\diag{\operatorname{diag}} \newcommand\Var{\operatorname{Var}} \newcommand\proj{\operatorname{proj}} \newcommand\pkg\textsf
\usepackage{bookmark}
\IfFileExists{xurl.sty}{\usepackage{xurl}}{} % add URL line breaks if available
\urlstyle{same}
% fallback for those not using the hyperref driver hyperxmp:
\makeatletter
\@ifundefined{xmpquote}{\newcommand{\xmpquote}[1]{#1}}{}
\makeatother
\hypersetup{
  pdftitle={Balancing constriction and expansion for domain-specific grammars},
  pdfauthor={Jason Cory Brunson},
  hidelinks,
  pdfcreator={LaTeX via pandoc}}

\title{Balancing constriction and expansion for domain-specific
grammars}
\author{Jason Cory Brunson}
\date{}

\begin{document}
\maketitle
\begin{abstract}
Graphical grammars distinguish from heuristic typologies in constituting
closed formal systems whose product diversity arises from recombination
rather than from \emph{de novo} redesign. In practice, the scope of an
implementation is constrained both by what options have been implemented
for each component (for example, the polar coordinatization, the grouped
count statistic, and the rectangle geometric element) and by what
combinations are enabled (these three yielding a pie graphic). In the
interest of flexibility, implementations are generally permissive of
syntactically valid combinations. However, within some domains,
universal or conditional rules may be in place to ensure that graphics
faithfully or efficiently communicate important properties of the input
data, such that grammatical implementations may compromise their
integrity by permitting violations of these rules. I will present
geometric data analysis as a motivating case study for this situation.
Visualizations of point clouds, and especially the use of biplots,
relies universally on a unitary aspect ratio and conditionally on data
sphering for plot elements to approximate geometric relationships among
data elements. Rather than a limitation of graphical grammar, these
should be viewed as opportunities to explore the phrase spaces of
different graphical domains. Indeed, the geometric constraints above
enable the meaningful use of two families of statistical
transformation---projections of constructions in higher dimensions and
constructions in relation to referent points---that would otherwise
produce uninterpretable graphics. This leads me to argue for a change in
language from Wilkinson's assertion that, while a grammatical system can
produce hideous graphics, it cannot produce meaningless ones. I suggest
that meaningfulness and meaninglessness are not mathematical properties
secured by formal rules but semantic properties vested with cultural
significance, and that graphical grammars must adapt to the culture of a
domain in order to serve as trustworthy tools for communicating
information within it.
\end{abstract}

\pagebreak

\subsection{Grammar of Graphics}\label{grammar-of-graphics}

In his book introducing the grammar of graphics (GoG), Wilkinson (2005)
emphasizes the value not only to the rigor but also to the extensibility
of taking a grammatical, in opposition to a typological, approach to
statistical graphics:

\begin{quote}
The remarkable consequence of building a closed formal system is that,
while it solves more complex applied problems, it can appear more
adaptive to the user. Paradoxically, closed systems often behave more
``openly'' than open systems.
\end{quote}

This flexibility assumes and relies on a permissiveness toward the
``phrase space'' of graphics that may be realized through syntactically
valid combinations. Indeed, Wilkinson asserts that, while his system can
produce many hideous graphics, it cannot produce meaningless ones. On my
reading, the primary focus, in the interest of meaningfulness, is on the
mathematical determinism connecting the data to the geometric elements
of a graphic---ultimately defined as a visual representation of the
abstract \emph{graph} of a mathematical function.

To be sure, Wilkinson details how meaningful graphics must also account
for human perceptual tendencies and for cultural norms. In popular
implementations, led by Wickham (2009), most directives toward good
practice take the form of defaults (e.g.~different color palettes for
different variable types) and warnings (e.g.~that opacity should not
encode a discrete variable). Wilkinson also posits certain
near-universal constraints, for example that the plotting window should
be sized to fit the input data, even if only summary statistics are
visualized.

What is missing in this account, i believe, is a reckoning with dialect:
Certain disciplines use specialized graphics whose uniquely structured
content imposes specialized constraints, so that extending the
grammatical approach to produce them requires constraining the
combinatorial freedom it makes possible. Rather than choosing between an
uninhibited extension that encourages misleading graphics and a
duplicative re-implementation for each domain, we should design and
adopt formal principles to link essential data properties to graphical
design patterns that must be sensitive to them.

\subsection{Geometric Data Analysis}\label{geometric-data-analysis}

GDA is a diverse family of techniques for analyzing data represented as
a \emph{point cloud}---usually, a finite multisubset
\(X \subset (\mathbb{R}^d, L^2)\)---and for transforming other data
types into point clouds for analysis. The points \(x_i \in X\) may be
attributed \emph{masses} \(m_i\), and the centroid of \(X\) is taken as
the origin, so that the point \emph{inertias}
\(m_i \lVert \vec{x} \rVert \vec{x}_i\) sum to zero. A foundational step
in GDA is the decomposition of inertia via singular value decomposition
(SVD): Taking \(X_{n \times m}\) to be centered and \(r \leq min(n,m)\),
we obtain
\[ X = U_{n \times r} \Sigma_{r \times r} {V_{m \times r}}^\top \] with
coordinatewise inertias
\(\operatorname{diag}(\Sigma) = (\sigma_1,\ldots,\sigma_r)\). The point
cloud \(F = U\Sigma\) is a rotation of an optimal approximation to
\(X\), with equality when \(r = \min(n,m)\); the column vectors
\(\vec{v}_1,\ldots,\vec{v}_r\) of \(V\) each maximize
\(\operatorname{Var}(\operatorname{proj}_{v_i}F_{i-1})\), where
\(F_{i-1}\) is the projection onto the previous \(v_{i-1}\) (starting
with \(F_0 = F\)); and each
\(\operatorname{Var}(\operatorname{proj}_{v_i}F_{i-1}) = {\sigma_i}^2\).
SVD thus obtains an ordered set of \emph{artificial coordinates} for the
proxy point cloud \(F\).

Two of the most widely used GDA techniques are principal components
analysis (PCA), which obtains \(X\) by centering and scaling numeric
case--variable data, and correspondence analysis (CA), which obtains a
matrix of standardized residuals from bivariate count data. These
techniques are called \emph{ordination} as opposed to regression (being
unsupervised rather than supervised) and \emph{gradient analysis} as
opposed to cluster analysis (yielding continuous rather than discrete
artificial variables). The artificial coordinates (\emph{principal
coordinates} along \emph{principal axes}) are not specified from a model
but derived from data; because they carry no intrinsic meaning, they are
sometimes called \emph{scaffolding}.

The SVD supports a unique statistical graphic, the \emph{biplot},
introduced by Gabriel (1971). Because the matrices \(U\) and \(V\) share
their \(r\) artificial column variables, row markers for both (proxy row
and column markers for \(X\)) can be overlaid on those axes. When \(F\)
and \(V\) are overlaid, the scalar product of the \(i^\text{th}\) row
marker and the \(j^\text{th}\) column marker approximates \(x_{ij}\);
ditto \(U\) and \(G = VD\). In general, a biplot overlay of \(UD^a\) and
\(VD^b\) is called \emph{balanced} when \(a + b = 1\), so that
row--column inner products approximate measurements. The inertia-laden
matrices \(F\) and \(G\) are said to contain \emph{principal
coordinates} for the rows and columns of \(X\), while \(U\) and \(V\)
contain \emph{standard coordinates} of unit magnitude. When a matrix
factor is represented in principal coordinates, the graphic approximates
its geometry in the standard coordinate space of the other factor: In
PCA, distances between the markers of \(F\) approximate those between
the rows of \(X\) in the space of measured variables, while inner
products and cosines between the markers of \(G\) approximate
covariances and correlations between the columns of \(X\) across the
case measurements. In CA, distances between markers of \(F\) and of
\(G\) approximate \(\chi^2\) distances between row and column profiles
in the count table.

Greenacre (2010) uses biplots to surface the mathematical unity of these
techniques. However, this unity is rarely reflected in software
packages. Implementations of PCA, CA, and several other techniques in
\textsf{stats} and \textsf{MASS}, for example, share no source code
beyond \texttt{base::svd()}, accept dissimilar arguments even for
equivalent parameters (e.g.~the rank \(r\)), and produce output classes
with few shared methods. My attempt to harmonize these tools through
\textsf{ordr}, a tidyverse-inspired API that foregrounds this common
theory, is the impetus for the remaining sections.

\subsection{Tensions between GDA and
GoG}\label{tensions-between-gda-and-gog}

Point clouds are more homogeneous and rigid mathematical abstractions
than the case--variable tables for which most statistical graphics are
designed. Most conspicuously, the coordinates must be equally scaled to
faithfully represent geometric relationships among the markers. This
need has not, to my knowledge, received focused attention in the GoG
space. For example, while Wilkinson showcases such graphics as Voronoi
polygons and sun glyphs that rely for their interpretations on a unit
aspect ratio, he only briefly mentions that this choice should be
informed by context. On the implementation side, \textsf{ggplot2}
provides the coordinatization ggproto \texttt{CoordFixed} with default
\texttt{ratio\ =\ 1}, but its invocation is the responsibility of the
user. This mismatch between domain and convention has enabled data
analysts to publish biplots that misrepresent the geometry of their
data.

Some default and hardcoded settings of \textsf{ggplot2} are also at odds
with best practices in GDA visualization: The artificial coordinates
used as scaffolding rarely warrant labels, let alone tick marks or grid
lines; instead, Gower et al. (2011) recommends that variable axes be
incorporated as graphical elements within the plot window.

Finally, as imbalanced biplots are strongly discouraged, rigorous
graphical software should inhibit such choices of the exponents \(a\)
and \(b\). Biplot-specific programs usually hide or protect this choice
from the user: \texttt{stats::biplot.princomp()} passes the argument
\texttt{scale} to \(b\) and \texttt{1\ -\ scale} to \(a\), while
\textsf{biplotEZ} provides the character-valued argument
\texttt{CA(variant)} and the logical-valued argument
\texttt{PCA(correlation.biplot)} that specify row-principal,
column-principal, or symmetric (neither-principal) overlays. Because
\textsf{ggplot2} and similar GoG implementations generally determine
plot layers from one data set at a time, this balance is again the
responsibility of the user.

\subsection{Contributions of GDA to
GoG}\label{contributions-of-gda-to-gog}

The aforementioned tensions should be taken not as intrinsic to the
grammatical approach but as an opportunity to advance theory and
practice. The \textsf{ggplot2} extension community in particular
maintains a tradition of embracing such challenges. I wish to highlight
three contributions i believe GDA has to make to GoG, with preliminary
implementations in \textsf{ordr} and its spinoff \textsf{gggda}.

First, it is often of interest in GDA to visually qualify how well the
geometry in the first artificial dimensions (those used to plot)
captures that of the full space.\footnote{This is quantified as a ratio
  of partial to total variances, called \emph{adequacy}.} One approach
is to assemble a minimal spanning tree in high dimensions and overlay
its projection (Jolliffe 2002). Such layers require that the statistical
transformation object have access to the full point cloud---inconsistent
if not impossible in GoG as originally conceived. Drawing from a
tradition of set-valued aesthetics tracing back to \textsf{ggmosaic},
\textsf{gggda} provides a mechanism to pass coordinate data of arbitrary
dimension for statistical transformation before projecting for geometric
construction.

Second, the unadjusted overlay of two marker sets can yield messy
biplots. A common adjustment for PCA biplots is to offset calibrated
variable axes from the dense core of case markers. This allows for
clearer projection, by the reader or even rendered in the plot, from
select markers onto select axes (yielding inner products). To try and
err at satisfying offsets is impractical, and algorithms are available
that work in most cases; but they require at least summary information
about secondary data---the cases---to compute graphical elements for
primary data---the variables. \textsf{gggda} implements a family of
\emph{referential stats}, built on the prototype \texttt{StatRef}, that
take a secondary data frame as an additional argument. While
straightforward, this trick couples with a feature of \textsf{ordr} that
preserves both matrix factors of an SVD-based technique so that each by
default produces layers with the other as referent (if invoked).

Finally, to evoke the common scale of the artificial axes, and to reveal
the difference in data range between the first and second axes, biplot
windows are often made square, rather than reshaped to fit the data.
\textsf{gggda} provides a coordinatization ggproto \texttt{CoordRect} to
achieve this for any desired ratio of window dimensions. This in turn
supports the in-development paneling ggproto \texttt{FacetOrd}, which
yields a pairs plot for artificial coordinates that follows the rigid
requirements above while permitting both fixed and varying panel sizes.
These conventions could be used to plot any set of real-valued data but
to my knowledge have not previously been introduced.

\subsection{A Case for Domain-Constrained
Grammar}\label{a-case-for-domain-constrained-grammar}

The case studies above lead me to argue that GoG, both in theory and in
practice, should take a deliberate approach to the constraint sets on
statistical graphics associated with specific domains. \textsf{ggplot2}
itself deals delicately and implicitly with Simple Features data, for
example, by transforming segments between geographic coordinates into
approximate meridians and tucking the coordinatization constructor
\texttt{coord\_sf()} into the geometric element constructor
\texttt{geom\_sf()}; but users may violate good practice by overriding
these behaviors. At an opposite extreme, \textsf{ggraph} and
\textsf{ggtern} are designed to operate entirely without native
\textsf{ggplot2} constructors. I suggest that a middle ground is
required to satisfy the needs of users in, and especially those
straddling, diverse specialties: From each specialty should be derived a
\emph{graphical dialect}, constituting both permitted objects and
constraints on their combination. Such dialects should be based on
standards of good practice and formalized mathematically so that they
may be unambiguously implemented. At a low level, an implementation
should then have the ability to declare and adopt a dialect, for a
single plot (perhaps by using a special initializer) or for a section of
code (perhaps by setting a global option), with the effect that plot
assembly informatively enables domain-specific features and imposes
domain-specific constraints (perhaps by code completions and
suggestions). At a high level, a family of well-defined dialects should
help organize a grand unified GoG at an important meso-scale between its
basic objects and its space of all feasible plot types.

\protect\phantomsection\label{refs}
\begin{CSLReferences}{1}{1}
\bibitem[\citeproctext]{ref-Gabriel1971}
Gabriel, K. R. 1971. {``The Biplot Graphic Display of Matrices with
Application to Principal Component Analysis.''} \emph{Biometrika} 58
(3): 453--67. \url{https://doi.org/10.2307/2334381}.

\bibitem[\citeproctext]{ref-Gower2011}
Gower, John C., Sugnet Lubbe, and Niël J. Le Roux. 2011.
\emph{Understanding Biplots}. John Wiley \& Sons, Ltd.
\url{https://www.wiley.com/go/biplots}.

\bibitem[\citeproctext]{ref-Greenacre2010}
Greenacre, Michael J. 2010. \emph{Biplots in Practice}. Fundación BBVA.
\url{https://www.fbbva.es/en/publicaciones/biplots-in-practice-7/}.

\bibitem[\citeproctext]{ref-Jolliffe2002}
Jolliffe, I. T. 2002. \emph{Principal Component Analysis}. 2nd ed.
Springer Series in Statistics. Springer-Verlag.
\url{https://doi.org/10.1007/b98835}.

\bibitem[\citeproctext]{ref-Wickham2009}
Wickham, Hadley. 2009. \emph{Ggplot2: Elegant Graphics for Data
Analysis}. Springer. \url{https://doi.org/10.1007/978-0-387-98141-3}.

\bibitem[\citeproctext]{ref-Wilkinson2005}
Wilkinson, Leland. 2005. \emph{The Grammar of Graphics}. 2nd ed.
Statistics and Computing. Springer-Verlag.
\url{https://doi.org/10.1007/0-387-28695-0}.

\end{CSLReferences}

\end{document}
